\documentclass[a4paper]{memoir}
\usepackage{graphicx}
%% Language and font encodings
\usepackage{microtype}
\usepackage[paperheight=88.9mm,paperwidth=63.5mm,margin=5mm]{geometry}
\usepackage{ragged2e}
\usepackage{adjustbox}
\usepackage{emoji}
\usepackage{pgffor}
\usepackage{tikz}
\newcommand{\cardtext}{0-10}
\newcommand{\grouptext}{2-4}
\newcommand{\agetext}{12+}
\newcommand{\clocktext}{20$''$}
\newcommand{\languagecrop}{french_crop}
\newcommand{\colorcount}{6}
\newcommand{\cardtitle}{}
\newcommand{\cardsubtitle}{}
\def\colors{0/myyellow, 1/myorange, 2/myred, 3/mypurple, 4/myblue, 5/mygreen}
\begin{document}
{\footnotesize
\input{../../../rule_lib/rule_lib.tex}
\noindent
\textbf{\emoji{party-popper} Variante Skyzone.}  
On garde les règles de Skyzone (ou Skyzone shape), mais il n’y a plus de défausse centrale.
\\
\\
\noindent
\textbf{Chaque joueur a sa propre défausse.}
\\
\\
\textbf{Quand tu pioches,} tu peux choisir entre la pioche centrale ou n’importe quelle défausse (y compris la tienne).
\\
\\
\textbf{Quand tu défausses,} tu poses ta carte sur ta propre défausse.
}
\newpage
\textbf{}
\end{document}

